\centerline{\bf \Huge Математическая абака}

\begin{flushright}
        \scriptsize{Если нет выхода - сделай его сам(ам-ам-ам).}
\end{flushright}

\section{Судоку}
В судоку необходимо заполнить пустые клетки цифрами от 1 до 9 (в первой задаче — фигурами 4 видов) так, чтобы в каждой строке и каждом столбце, а также в каждой отделённой жирными границами зоне было ровно по одной цифре (фигуре) каждого вида.

\problem{10}
В этой задаче вместо цифр используются фигуры 4 видов (четвёртый вид фигур можно придумать самостоятельно, так как он не задан в условии).
\begin{center}
\includegraphics[width=0.3\textwidth]{Летние Сборы/Сборы Элиста 2024/ВВ/судоку_1.png}
\end{center}

\problem{20}
\begin{center}
\includegraphics[width=0.2\textwidth]{Летние Сборы/Сборы Элиста 2024/ВВ/судоку_2.png}
\end{center}

\problem{30}
\begin{center}
\includegraphics[width=0.2\textwidth]{Летние Сборы/Сборы Элиста 2024/ВВ/судоку_3.png}
\end{center}

\newpage

\problem{40}
\begin{center}
\includegraphics[width=0.4\textwidth]{Летние Сборы/Сборы Элиста 2024/ВВ/судоку_4.png}
\end{center}

\problem{50}
\begin{center}
\includegraphics[width=0.4\textwidth]{Летние Сборы/Сборы Элиста 2024/ВВ/судоку_5.png}
\end{center}

\problem{60}
\begin{center}
\includegraphics[width=0.4\textwidth]{Летние Сборы/Сборы Элиста 2024/ВВ/судоку_6.png}
\end{center}

\newpage

\section{Расстановка скобок и знаков}

\problem{10}
Расставьте знаки действий, где это необходимо, чтобы равенство стало верным: 2 0 1 2 = 18

\problem{20}
Расставьте в снежинке все цифры от 1 до 7 так, чтобы все суммы по три числа на отрезках были одинаковыми:

\problem{30}
Решите ребус (одинаковым буквам соответствуют одинаковые цифры, разным - разные):
\begin{center}

\includegraphics[width=0.3\textwidth]{Летние Сборы/Сборы Элиста 2024/ВВ/скобки_3.png}
\end{center}

\problem{40}
Расставьте числа от 1 до 12 в кружки так, чтобы каждая стрелка вела от большего числа к меньшему.
\begin{center}
\includegraphics[width=0.2\textwidth]{Летние Сборы/Сборы Элиста 2024/ВВ/скобки_4.png}
\end{center}

\problem{50}Расставьте числа в пропуски так, чтобы сумма любых четырёх чисел, идущих подряд, равнялась 29: 
$$15 \_ \_ \_ \_ \_ \_ 7 \_ 3$$

\problem{60}
Вася написал на доске 9 троек в ряд и перед каждой цифрой, включая первую, поставил знак «+» или «–», а потом расставил пары скобок, сколько ему вздумалось. Петя подошёл и посчитал значение получившегося выражения. Сколько возможных ответов он мог получить?

\newpage

\section{Логика}



\problem{10}
Илюша получил пятерку по литературе и истории. Аннушка сказала: "Илюша получил пятерку по истории или по биологии". Уж не соврала ли Аннушка?

\problem{20}
Правда ли, что в этом помещении все драконы добрые?

\problem{30}
Ниф-ниф, Наф-наф и Нуф-нуф пришли в казино в серебряной, золотой и платиновой цепях. Перстни у них были из тех же металлов. У Ниф-нифа цепь и перстень были изготовлены из одинакового металла. У Наф-нафа ни цепь, ни перстень не были серебряными. У Нуф-нуфа цепь была из платины, а перстень не из платины. У кого из чего были изготовлены украшения?

\problem{40}
Следующая беседа имела место между чиновником Министерства социального обеспечения одной далёкой страны и претендентом на детское пособие:

— Сколько у вас детей?

— Трое.

— Сколько им лет?

— Произведение их возрастов равно 36.

— Вы должны дать мне больше информации.

— Сумма их возрастов равняется номеру магазина напротив вашего офиса.

— Этого ещё не достаточно.

— Вы правы. Могу ещё сказать, что старший сын рыжий.

— Теперь мне всё понятно...

Сколько лет детям?


\problem{50}
За круглым столом сидят 10 человек. Требуется установить, сколько рыцарей может быть за столом, если каждый из сидящих произносит одну и ту же фразу: «Оба моих соседа — лжецы».


\problem{60}
Из 16 монет одна — фальшивая, причем неизвестно, легче она настоящих или тяжелее. Найти ее за 4 взвешивания.
